% Generated draft: requires mathematical and source review.
\documentclass[11pt]{article}
\usepackage[margin=25mm]{geometry}
\usepackage{fontspec}
\setmainfont{Latin Modern Roman}
\usepackage{amsmath,amssymb,mathtools,unicode-math}
\setmathfont{Latin Modern Math}
\usepackage{enumitem,needspace,xurl,hyperref}
\setlength{\parindent}{0pt}
\setlength{\parskip}{5pt}
\setlength{\emergencystretch}{3em}
\newcommand{\N}{\mathbb{N}}
\newcommand{\Z}{\mathbb{Z}}
\newcommand{\Q}{\mathbb{Q}}
\newcommand{\R}{\mathbb{R}}
\newcommand{\C}{\mathbb{C}}
\newcommand{\F}{\mathbb{F}}
\newcommand{\A}{\mathbb{A}}
\newcommand{\PP}{\mathbb{P}}
\newcommand{\E}{\mathbb{E}}
\newcommand{\eps}{\varepsilon}
\newcommand{\dd}{\,\mathrm{d}}
\begin{document}
\section*{UnsolvedMath ID: 800694}
\textbf{Source problem code:} catalogue entry
\subsection{Definitions and mathematical statement}
In the central simple algebra setting of the supplied source, interpret an \(F\)-division algebra as a finite-dimensional central division algebra \(A\) over a field \(F\). Write \(n=\deg(A)\), so that \(\dim_F A=n^2\). The supplied statement does not explicitly state centrality or finite dimensionality; this interpretation requires review.

A separable maximal subfield of \(A\) is a subfield \(K\subseteq A\) containing \(F\) such that \(K/F\) is separable and \([K:F]=n\). After choosing an \(F\)-embedding of \(K\) into a separable closure of \(F\), let \(L_K/F\) be its Galois closure and put
\[
G_K=\operatorname{Gal}(L_K/F).
\]
The extension \(L_K\) splits \(A\), meaning that
\[
A\otimes_F L_K\cong M_n(L_K).
\]
Thus the groups under consideration are those arising from Galois closures of separable maximal subfields, rather than from arbitrary splitting extensions. Define their isomorphism classes by
\[
\mathcal G(A)=\bigl\{[\operatorname{Gal}(L_K/F)]_{\cong}: K\subseteq A\text{ is a separable maximal subfield}\bigr\}.
\]
Each such group acts faithfully and transitively on the \(n\) distinct \(F\)-embeddings of \(K\) into \(L_K\), and hence has a permutation representation as a subgroup of \(S_n\).

\textbf{Problem.} Determine the smallest group or groups in \(\mathcal G(A)\). The source does not specify whether ``smallest'' refers to group order, subgroup containment in a permutation representation, or another comparison, so this part does not yet specify a unique minimization problem. Its initial reference to a maximal subfield also does not explicitly impose separability; the definition above uses separable maximal subfields to give the Galois-closure construction its usual meaning.

In particular, writing \(A_n\) for the alternating group on \(n\) letters, determine whether the following assertion holds:
\[
\begin{gathered}
\text{For every field }F,\text{ every integer }n\geq 1,\text{ and every finite-dimensional}\\
\text{central division }F\text{-algebra }A\text{ of degree }n,\text{ there exists a separable}\\
\text{maximal subfield }K\subseteq A\text{ such that }
\operatorname{Gal}(L_K/F)\cong A_n.
\end{gathered}
\]
No restriction on the characteristic of \(F\) or lower bound on the degree is supplied. The displayed universal formulation retains that unrestricted wording; the intended range of degrees, including boundary cases, requires clarification. This is a statement-only draft, not an assertion that the proposed universal statement is true.
\subsection{Short English statement}
For a central division algebra, consider the Galois groups obtained by taking Galois closures of its separable maximal subfields. Determine which of these groups are smallest, with the meaning of ``smallest'' still unspecified. In particular, ask whether every algebra of degree \(n\) admits such a subfield with Galois-closure group \(A_n\). The intended degree range and implicit algebra hypotheses require review.
\subsection{Sources}
\begin{itemize}
\item Local UnsolvedMath registry, supplied complete source record: permanent ID \(800694\), source problem code \texttt{catalogue entry}, title \emph{Open Problems on Central Simple Algebras --- Problem 4}. The supplied registry labels the problem \texttt{open}; this is a supplied status, not an independently verified mathematical or literature assessment.
\item Auel, Brussel, Garibaldi, and Vishne, \emph{Open Problems on Central Simple Algebras} (2010), Problem 4. Supplied source URL: \url{https://arxiv.org/abs/1006.3304}. The record attributes the problem to these authors and gives proposed year 2010.
\item Supplied provenance records access on 2026-10-16 and extraction by \texttt{arxiv\_latex}. These are supplied provenance claims; the source has not been independently accessed or verified for this draft. Supplied upstream revision: \texttt{d371ee917ac9647fbb06cd6fc529deec6f3dda82}.
\end{itemize}
\end{document}
